\documentclass[10pt]{article}

\usepackage[utf8]{inputenc}
\usepackage[T1]{fontenc}
\usepackage{amsmath,amssymb,amsthm}
\usepackage{booktabs}
\usepackage{enumitem}
\usepackage{geometry}
\usepackage{hyperref}
\usepackage[most]{tcolorbox}

\geometry{margin=0.75in}
\hypersetup{colorlinks=true,linkcolor=blue,urlcolor=blue}

\setlist{itemsep=1pt,topsep=2pt,parsep=0pt,partopsep=0pt}
\setlength{\abovedisplayskip}{5pt plus 2pt minus 2pt}
\setlength{\belowdisplayskip}{5pt plus 2pt minus 2pt}
\setlength{\abovedisplayshortskip}{3pt plus 2pt minus 1pt}
\setlength{\belowdisplayshortskip}{3pt plus 2pt minus 1pt}
\allowdisplaybreaks

\makeatletter
\renewcommand\section{\@startsection{section}{1}{\z@}%
  {-1.5ex \@plus -.4ex \@minus -.2ex}%
  {.7ex \@plus .2ex}%
  {\normalfont\large\bfseries}}
\renewcommand\subsection{\@startsection{subsection}{2}{\z@}%
  {-1.0ex \@plus -.3ex \@minus -.2ex}%
  {.35ex \@plus .1ex}%
  {\normalfont\normalsize\bfseries}}
\renewcommand\subsubsection{\@startsection{subsubsection}{3}{\z@}%
  {-.8ex \@plus -.2ex \@minus -.1ex}%
  {.25ex \@plus .1ex}%
  {\normalfont\normalsize\bfseries\itshape}}
\makeatother

\newtheoremstyle{compactdefinition}
  {4pt}{4pt}{\normalfont}{}{\bfseries}{.}{.5em}{}
\newtheoremstyle{compactproposition}
  {4pt}{4pt}{\normalfont}{}{\bfseries}{.}{.5em}{}
\theoremstyle{compactdefinition}
\newtheorem{definition}{Definition}[section]
\theoremstyle{compactproposition}
\newtheorem{proposition}{Proposition}[section]

\definecolor{PropositionBlue}{RGB}{31,78,121}
\tcolorboxenvironment{proposition}{
  enhanced,breakable,
  colback=PropositionBlue!4,colframe=PropositionBlue,
  boxrule=0.9pt,arc=1.5mm,
  left=2mm,right=2mm,top=1.5mm,bottom=1.5mm
}

\newcommand{\R}{\mathbb{R}}
\newcommand{\M}{\mathcal{M}}
\newcommand{\rank}{\operatorname{rank}}
\newcommand{\im}{\operatorname{Im}}
\newcommand{\tr}{\operatorname{tr}}
\newcommand{\diag}{\operatorname{diag}}
\newcommand{\Span}{\operatorname{span}}
\newcommand{\Var}{\operatorname{Var}}

\title{Summer Bootcamp -- Session 1 Notes\\Linear Algebra}
\author{Nathan Sauldubois}
\date{August 2026}

\begin{document}
\maketitle

\section*{Purpose of the Session}

This first three-hour session is a refresher on the linear algebra needed later in
the bootcamp. The goal is not to cover linear algebra exhaustively, but to make
students operational with the matrix language used in financial models,
statistics, optimization, and numerical computation.

These notes are more detailed than the slides. They are organized by conceptual
blocks, with definitions, propositions, and short proofs whenever a result is
used later:
\begin{itemize}
\item algebraic operations;
\item associated scalar and matrix operators;
\item inversion and linear systems;
\item injectivity, surjectivity, image, kernel, and rank;
\item classical decompositions, including diagonalization and Cholesky.
\end{itemize}

\section{Algebraic Operations}

\begin{definition}[Matrix and matrix-vector map]
A matrix is a rectangular array of real numbers,
\[
A=(a_{ij})_{1\le i\le n,\ 1\le j\le m}\in \M_{n,m}(\R).
\]
The convention used here is that $n$ is the number of rows and $m$ is the number
of columns. A matrix in $\M_{n,m}(\R)$ maps vectors in $\R^m$ to vectors in
$\R^n$.
\end{definition}

\subsection*{Numerical example}

\[
A=
\begin{pmatrix}
1&2&3\\
4&5&6
\end{pmatrix}
\in \M_{2,3}(\R),
\qquad
x=
\begin{pmatrix}
7\\8\\9
\end{pmatrix}
\in \R^3.
\]
The product $Ax$ is defined because $A$ has three columns and $x$ has three
entries. The result belongs to $\R^2$:
\[
Ax=
\begin{pmatrix}
1\cdot7+2\cdot8+3\cdot9\\
4\cdot7+5\cdot8+6\cdot9
\end{pmatrix}
=
\begin{pmatrix}
50\\122
\end{pmatrix}.
\]

\begin{definition}[Addition and scalar multiplication]
For two matrices of the same size,
\[
(A+B)_{ij}=a_{ij}+b_{ij},
\qquad
(\lambda A)_{ij}=\lambda a_{ij}.
\]
Addition and scalar multiplication are entrywise operations.
\end{definition}

\begin{proposition}[Entrywise operations behave linearly]
For matrices of the same size and scalars $\lambda,\mu\in\R$,
\[
A+B=B+A,\qquad (A+B)+C=A+(B+C),
\]
\[
\lambda(A+B)=\lambda A+\lambda B,\qquad
(\lambda+\mu)A=\lambda A+\mu A.
\]
\end{proposition}

\begin{proof}
All identities are checked entry by entry. For instance,
\[
(\lambda(A+B))_{ij}=\lambda(a_{ij}+b_{ij})
=\lambda a_{ij}+\lambda b_{ij}
=(\lambda A+\lambda B)_{ij}.
\]
The other identities follow from the same identities in $\R$.
\end{proof}

\subsection*{Example: addition and scalar multiplication}

Let
\[
A=\begin{pmatrix}1&-2&3\\0&4&-1\end{pmatrix},
\qquad
B=\begin{pmatrix}2&0&-1\\-3&1&5\end{pmatrix}.
\]
Then
\[
A+B=
\begin{pmatrix}
3&-2&2\\
-3&5&4
\end{pmatrix},
\]
and
\[
A-2B=
\begin{pmatrix}
1-4 & -2-0 & 3-(-2)\\
0-(-6) & 4-2 & -1-10
\end{pmatrix}
=
\begin{pmatrix}
-3&-2&5\\
6&2&-11
\end{pmatrix}.
\]

\subsection*{Example: a linear combination of vectors}

Let
\[
u=\begin{pmatrix}1\\2\\0\end{pmatrix},
\qquad
v=\begin{pmatrix}-1\\1\\3\end{pmatrix}.
\]
Then
\[
3u-2v
=
3\begin{pmatrix}1\\2\\0\end{pmatrix}
-2\begin{pmatrix}-1\\1\\3\end{pmatrix}
=
\begin{pmatrix}5\\4\\-6\end{pmatrix}.
\]
This is the same vector-space operation used with the columns of a matrix.

\begin{definition}[Matrix multiplication]
If $A\in \M_{n,m}(\R)$ and $B\in \M_{m,p}(\R)$, then $AB$ is defined and belongs
to $\M_{n,p}(\R)$. Its entries are
\[
(AB)_{ij}=\sum_{k=1}^{m}a_{ik}b_{kj}.
\]
The number of columns of $A$ must equal the number of rows of $B$.
\end{definition}

\subsection*{Dimension checks}

Students should constantly track dimensions. Most beginner mistakes with matrix
algebra are dimension mistakes:
\[
(2\times3)(3\times1)\to 2\times1,
\qquad
(4\times2)(2\times5)\to 4\times5,
\]
whereas
\[
(2\times3)(2\times5)
\]
is not defined.

\begin{proposition}[Matrix multiplication represents composition]
Matrix multiplication represents composition of linear transformations. If
$B$ first maps a vector from $\R^p$ to $\R^m$, and $A$ then maps from $\R^m$ to
$\R^n$, the composed transformation is $AB:\R^p\to\R^n$.
\end{proposition}

\begin{proof}
For $x\in\R^p$, the $i$th coordinate of $(AB)x$ is
\[
\sum_{j=1}^p(AB)_{ij}x_j
=\sum_{j=1}^p\sum_{k=1}^m a_{ik}b_{kj}x_j
=\sum_{k=1}^m a_{ik}\left(\sum_{j=1}^p b_{kj}x_j\right).
\]
This is exactly the $i$th coordinate of $A(Bx)$.
\end{proof}

\subsection*{Worked example}

Let
\[
A=\begin{pmatrix}1&2&-1\\0&3&4\end{pmatrix},
\qquad
B=\begin{pmatrix}2&1\\-1&0\\3&2\end{pmatrix}.
\]
Then
\[
AB=
\begin{pmatrix}
1\cdot2+2(-1)+(-1)3 & 1\cdot1+2\cdot0+(-1)2\\
0\cdot2+3(-1)+4\cdot3 & 0\cdot1+3\cdot0+4\cdot2
\end{pmatrix}
=
\begin{pmatrix}-3&-1\\9&8\end{pmatrix}.
\]

\subsection*{Matrix-vector product as a linear combination of columns}

Let
\[
A=
\begin{pmatrix}
2&-1&3\\
0&4&1
\end{pmatrix},
\qquad
x=
\begin{pmatrix}
1\\-2\\0
\end{pmatrix}.
\]
Then
\[
Ax=
1\begin{pmatrix}2\\0\end{pmatrix}
-2\begin{pmatrix}-1\\4\end{pmatrix}
+0\begin{pmatrix}3\\1\end{pmatrix}
=
\begin{pmatrix}4\\-8\end{pmatrix}.
\]
Thus $Ax$ is a linear combination of the columns of $A$, with coefficients given
by the entries of $x$.

\begin{proposition}[Product rules]
Matrix multiplication is associative and distributive whenever the dimensions
are compatible. It is not commutative in general.
\end{proposition}

\begin{proof}
Associativity follows from the composition interpretation:
$(AB)C$ and $A(BC)$ both send $x$ to $A(B(Cx))$. Distributivity follows from
linearity of matrix-vector maps, since both $A(B+C)$ and $AB+AC$ send $x$ to
$A(Bx+Cx)$. The example below proves that commutativity fails in general.
\end{proof}

In general, $AB\neq BA$. For
\[
A=\begin{pmatrix}1&2\\0&1\end{pmatrix},
\qquad
B=\begin{pmatrix}0&1\\1&0\end{pmatrix},
\]
we get
\[
AB=\begin{pmatrix}2&1\\1&0\end{pmatrix},
\qquad
BA=\begin{pmatrix}0&1\\1&2\end{pmatrix}.
\]
Thus this single example is enough to show that matrix multiplication is not
commutative in general.

\section{Associated Operators: Transpose, Trace, and Determinant}

\begin{definition}[Transpose]
The transpose of $A=(a_{ij})$ is the matrix $A^\top$ satisfying
\[
(A^\top)_{ij}=a_{ji}.
\]
It turns rows into columns and columns into rows.
\end{definition}

For example,
\[
A=
\begin{pmatrix}
1&3&-2\\
0&4&5
\end{pmatrix}
\quad\Longrightarrow\quad
A^\top=
\begin{pmatrix}
1&0\\
3&4\\
-2&5
\end{pmatrix}.
\]

\begin{proposition}[Transpose rules]
Whenever the dimensions are compatible,
\[
(A+B)^\top=A^\top+B^\top,\qquad
(\lambda A)^\top=\lambda A^\top,\qquad
(AB)^\top=B^\top A^\top.
\]
For vectors $x,y$, the dot product can be written as $x^\top y$, and
\[
(Ax)^\top y=x^\top A^\top y.
\]
\end{proposition}

\begin{proof}
The first two identities are entrywise. For the product,
\[
\bigl((AB)^\top\bigr)_{ij}=(AB)_{ji}
=\sum_k a_{jk}b_{ki}
=\sum_k (B^\top)_{ik}(A^\top)_{kj}
=(B^\top A^\top)_{ij}.
\]
The dot-product identity follows by applying this product rule to $(Ax)^\top y$.
\end{proof}

\begin{definition}[Symmetric and positive semidefinite matrices]
A square matrix $A$ is symmetric if $A^\top=A$. It is positive semidefinite if
\[
x^\top A x\ge 0
\qquad\text{for all }x\in\R^n.
\]
It is positive definite if the inequality is strict for every nonzero $x$.
\end{definition}

\begin{proposition}[$X^\top X$ is positive semidefinite]
For every real matrix $X$, the matrix $X^\top X$ is symmetric and positive
semidefinite.
\end{proposition}

\begin{proof}
Symmetry follows from $(X^\top X)^\top=X^\top(X^\top)^\top=X^\top X$. For any
vector $z$,
\[
z^\top X^\top Xz=(Xz)^\top(Xz)=\|Xz\|^2\ge0.
\]
Thus $X^\top X$ is positive semidefinite.
\end{proof}

For
\[
A=
\begin{pmatrix}
2&1\\
1&2
\end{pmatrix},
\]
we have
\[
x^\top Ax=2x_1^2+2x_1x_2+2x_2^2
=(x_1+x_2)^2+x_1^2+x_2^2.
\]
This is strictly positive for $x\neq0$, so $A$ is positive definite.

\begin{definition}[Trace]
For a square matrix $A$, the trace is the sum of diagonal entries:
\[
\tr(A)=\sum_{i=1}^n a_{ii}.
\]
\end{definition}

For example,
\[
\tr
\begin{pmatrix}
3&1&0\\
2&-4&5\\
1&1&7
\end{pmatrix}
=3-4+7=6.
\]

\begin{proposition}[Trace rules]
Whenever the expressions are defined,
\[
\tr(A+B)=\tr(A)+\tr(B),
\qquad
\tr(\lambda A)=\lambda\tr(A),
\]
\[
\tr(AB)=\tr(BA)
\]
whenever both products make sense as square matrices.
\end{proposition}

\begin{proof}
Linearity is entrywise on the diagonal. If $A\in\M_{n,m}(\R)$ and
$B\in\M_{m,n}(\R)$, then
\[
\tr(AB)=\sum_{i=1}^n\sum_{k=1}^m a_{ik}b_{ki}
=\sum_{k=1}^m\sum_{i=1}^n b_{ki}a_{ik}
=\tr(BA).
\]
\end{proof}

\begin{definition}[Determinant]
The determinant is a scalar attached to a square matrix. It measures signed
volume scaling and detects whether a square matrix loses dimension.
\end{definition}

In $\R^2$, $|\det(A)|$ is the area scaling factor of the linear map $x\mapsto
Ax$. In $\R^3$, $|\det(A)|$ is the volume scaling factor. If $\det(A)=0$, some
volume collapses to zero, so the transformation cannot be invertible.

\begin{definition}[$2\times2$ determinant]
For
\[
A=
\begin{pmatrix}
a&b\\
c&d
\end{pmatrix},
\]
the determinant is
\[
\det(A)=ad-bc.
\]
\end{definition}

\subsection*{Numerical examples}

\[
\det\begin{pmatrix}2&3\\1&-4\end{pmatrix}
=2(-4)-3(1)=-11.
\]
The determinant is nonzero, so the matrix is invertible.

\[
\det\begin{pmatrix}1&2\\2&4\end{pmatrix}
=1\cdot4-2\cdot2=0.
\]
The determinant is zero because the second row is twice the first. The rows do
not span $\R^2$, and the matrix is not invertible.

\begin{definition}[General determinant: the Leibniz formula]
Let $A=(a_{ij})\in\M_n(\R)$ and let $\mathfrak S_n$ denote the set of all
permutations of $\{1,\ldots,n\}$. The determinant is
\[
\boxed{
\det(A)
=\sum_{\sigma\in\mathfrak S_n}
\operatorname{sgn}(\sigma)
\prod_{i=1}^n a_{i,\sigma(i)}
}.
\]
Here $\operatorname{sgn}(\sigma)=1$ for an even permutation and
$\operatorname{sgn}(\sigma)=-1$ for an odd permutation. Equivalently,
\[
\operatorname{sgn}(\sigma)=(-1)^{N(\sigma)},
\qquad
N(\sigma)=\#\{(i,j):i<j,\ \sigma(i)>\sigma(j)\},
\]
where $N(\sigma)$ is the number of inversions of $\sigma$.
\end{definition}

Each term selects exactly one entry from every row and every column. Since
$\mathfrak S_n$ contains $n!$ permutations, the formula has $n!$ signed terms.
It is the clean general definition but a deliberately horrible computational
method for large $n$; elimination or a matrix factorization is far more
efficient.

\begin{definition}[Minors and cofactors]
For $A\in\M_n(\R)$, let $A_{ij}$ be the $(n-1)\times(n-1)$ matrix obtained by
deleting row $i$ and column $j$. The corresponding minor and cofactor are
\[
M_{ij}=\det(A_{ij}),
\qquad
C_{ij}=(-1)^{i+j}M_{ij}.
\]
The signs of the cofactors follow the checkerboard pattern
\[
\begin{pmatrix}
+&-&+&\cdots\\
-&+&-&\cdots\\
+&-&+&\cdots\\
\vdots&\vdots&\vdots&\ddots
\end{pmatrix}.
\]
\end{definition}

\begin{proposition}[Laplace expansion along a row or a column]
For any fixed row $i$,
\[
\boxed{\det(A)=\sum_{j=1}^n a_{ij}C_{ij}
=\sum_{j=1}^n(-1)^{i+j}a_{ij}\det(A_{ij}).}
\]
For any fixed column $j$,
\[
\boxed{\det(A)=\sum_{i=1}^n a_{ij}C_{ij}
=\sum_{i=1}^n(-1)^{i+j}a_{ij}\det(A_{ij}).}
\]
Thus one may expand along any row or any column. In practice, choose one that
contains many zeros.
\end{proposition}

\subsection*{The $3\times3$ determinant}

For
\[
A=
\begin{pmatrix}
a&b&c\\
d&e&f\\
g&h&i
\end{pmatrix},
\]
expansion along the first row gives
\[
\det(A)
=
a
\begin{vmatrix}
e&f\\
h&i
\end{vmatrix}
-
b
\begin{vmatrix}
d&f\\
g&i
\end{vmatrix}
+
c
\begin{vmatrix}
d&e\\
g&h
\end{vmatrix}.
\]
Equivalently,
\[
\det(A)=a(ei-fh)-b(di-fg)+c(dh-eg).
\]

\subsection*{Numerical $3\times3$ example}

Let
\[
A=
\begin{pmatrix}
1&2&0\\
-1&3&1\\
2&0&4
\end{pmatrix}.
\]
Expanding along the first row,
\[
\det(A)
=
1
\begin{vmatrix}
3&1\\
0&4
\end{vmatrix}
-
2
\begin{vmatrix}
-1&1\\
2&4
\end{vmatrix}
+0
\begin{vmatrix}
-1&3\\
2&0
\end{vmatrix}.
\]
Thus
\[
\det(A)=1(12)-2((-1)4-1\cdot2)=12-2(-6)=24.
\]
Since $\det(A)\neq0$, this matrix is invertible.

\begin{proposition}[Triangular determinant]
If $A$ is upper or lower triangular, then
\[
\det(A)=\prod_{i=1}^n a_{ii}.
\]
\end{proposition}

\begin{proof}
In the Leibniz formula for the determinant, every non-identity permutation
chooses at least one entry below the diagonal for an upper triangular matrix
(or above the diagonal for a lower triangular matrix). That entry is zero, so
only the diagonal product remains.
\end{proof}

For example,
\[
\det
\begin{pmatrix}
2&-1&4\\
0&3&5\\
0&0&-2
\end{pmatrix}
=2\cdot3\cdot(-2)=-12.
\]

\begin{proposition}[Determinant rules]
For square matrices of the same size:
\begin{itemize}
\item $\det(I_n)=1$.
\item $\det(AB)=\det(A)\det(B)$.
\item $\det(A^\top)=\det(A)$.
\item If $A$ is invertible, then $\det(A^{-1})=1/\det(A)$.
\item Swapping two rows changes the sign of the determinant.
\item Multiplying one row by $\lambda$ multiplies the determinant by $\lambda$.
\item Adding a multiple of one row to another row does not change the determinant.
\end{itemize}
\end{proposition}

\begin{proof}
The row-operation rules come from the determinant being multilinear in the rows
and alternating: two equal rows force determinant zero, so swapping rows changes
sign and adding a multiple of one row to another contributes only a determinant
with two proportional rows. The identity and transpose rules follow from the
Leibniz formula. The product rule $\det(AB)=\det(A)\det(B)$ follows because
applying $B$ and then $A$ multiplies oriented volume first by $\det(B)$ and then
by $\det(A)$. Once the product rule is known,
$1=\det(I_n)=\det(AA^{-1})=\det(A)\det(A^{-1})$, which gives the inverse rule.
\end{proof}

\section{Inversion and Linear Systems}

A linear system can be written as $Ax=b$. The matrix $A$ contains coefficients,
$x$ contains unknowns, and $b$ is the target vector.

\begin{definition}[Identity and inverse]
The identity matrix $I_n$ satisfies $AI_n=A$ and $I_nA=A$ whenever dimensions are
compatible. A square matrix $A$ is invertible if there exists $A^{-1}$ such that
\[
AA^{-1}=A^{-1}A=I_n.
\]
\end{definition}

\begin{proposition}[$2\times2$ invertibility and inverse formula]
For
\[
A=\begin{pmatrix}a&b\\c&d\end{pmatrix},
\]
the matrix is invertible if and only if $\det(A)=ad-bc\neq0$. In that case,
\[
\boxed{
A^{-1}=\frac{1}{\det(A)}
\begin{pmatrix}d&-b\\-c&a\end{pmatrix}
=\frac{1}{ad-bc}
\begin{pmatrix}d&-b\\-c&a\end{pmatrix}.}
\]
\end{proposition}

\begin{proof}
If $ad-bc\neq0$, direct multiplication of $A$ by the displayed matrix gives
$I_2$ on both sides. Conversely, if $ad-bc=0$, the rows are linearly dependent,
so a nonzero vector belongs to $\ker(A)$. Such a matrix cannot be invertible.
\end{proof}

\begin{proposition}[Invertible systems have one solution]
If $A$ is invertible, the system $Ax=b$ has the unique solution
\[
x=A^{-1}b.
\]
\end{proposition}

\begin{proof}
Multiplying $Ax=b$ by $A^{-1}$ gives $x=A^{-1}b$, so any solution must be this
vector. Conversely, $A(A^{-1}b)=b$, so this vector is indeed a solution.
\end{proof}

For numerical work, however, solving $Ax=b$ directly is usually preferable to
forming $A^{-1}$ explicitly.

\subsection*{Example: solving a $2\times2$ system}

Consider
\[
\begin{cases}
2x_1+3x_2=5,\\
x_1-4x_2=-2.
\end{cases}
\]
Then
\[
A=\begin{pmatrix}2&3\\1&-4\end{pmatrix},
\qquad
b=\begin{pmatrix}5\\-2\end{pmatrix}.
\]
The determinant is
\[
\det(A)=2(-4)-3(1)=-11\neq0,
\]
so there is a unique solution. From the second equation, $x_1=-2+4x_2$.
Substituting into the first equation gives
\[
2(-2+4x_2)+3x_2=5
\quad\Rightarrow\quad
11x_2=9
\quad\Rightarrow\quad
x_2=\frac{9}{11}.
\]
Then
\[
x_1=-2+4\frac{9}{11}=\frac{14}{11}.
\]
Thus
\[
x=
\begin{pmatrix}
14/11\\
9/11
\end{pmatrix}.
\]

\subsection*{Solving using the inverse}

For a $2\times2$ matrix, the inverse formula gives
\[
A^{-1}
=
\frac{1}{-11}
\begin{pmatrix}
-4&-3\\
-1&2
\end{pmatrix}
=
\begin{pmatrix}
4/11&3/11\\
1/11&-2/11
\end{pmatrix}.
\]
Then
\[
x=A^{-1}b
=
\begin{pmatrix}
4/11&3/11\\
1/11&-2/11
\end{pmatrix}
\begin{pmatrix}5\\-2\end{pmatrix}
=
\begin{pmatrix}
14/11\\
9/11
\end{pmatrix}.
\]

\begin{proposition}[Forward substitution]
If $L$ is lower triangular with nonzero diagonal entries, the system $Lx=b$ can
be solved uniquely from the first equation to the last. Similarly, an upper
triangular system is solved by backward substitution.
\end{proposition}

\begin{proof}
The first equation contains only $x_1$ and has a unique solution because
$\ell_{11}\neq0$. Once $x_1,\ldots,x_{i-1}$ are known, the $i$th equation has
only one unknown left, namely $x_i$, and $\ell_{ii}\neq0$ determines it
uniquely. This induction solves all coordinates.
\end{proof}

Triangular systems are solved sequentially. In fixed income, this is the linear
algebra behind basic bootstrapping.

Suppose
\[
P_1=0.97,\qquad P_2=0.92,\qquad P_3=0.855.
\]
Write
\[
D(T_1)=y_1,\qquad D(T_2)=y_1+y_2,\qquad D(T_3)=y_1+y_2+y_3.
\]
Then
\[
\begin{pmatrix}
1&0&0\\
1&1&0\\
1&1&1
\end{pmatrix}
\begin{pmatrix}y_1\\y_2\\y_3\end{pmatrix}
=
\begin{pmatrix}0.97\\0.92\\0.855\end{pmatrix}.
\]
Forward substitution gives
\[
y_1=0.97,\qquad y_2=0.92-0.97=-0.05,\qquad
y_3=0.855-0.92=-0.065.
\]

\section{Injectivity, Surjectivity, Image, Kernel, and Rank}

\begin{definition}[Linear map associated with a matrix]
A matrix $A\in\M_{m,n}(\R)$ defines a linear map
\[
T_A:\R^n\to\R^m,
\qquad
T_A(x)=Ax.
\]
The notions of injectivity and surjectivity are properties of this map.
\end{definition}

\begin{definition}[Image and surjectivity]
The image, or range, of $A$ is
\[
\im(A)=\{Ax:x\in\R^n\}\subset\R^m.
\]
The map $T_A:\R^n\to\R^m$ is surjective if every vector in $\R^m$ is reached:
\[
\im(A)=\R^m.
\]
\end{definition}

\begin{proposition}[Column span and surjectivity criterion]
The image of $A$ is the span of its columns. Therefore $T_A:\R^n\to\R^m$ is
surjective if and only if $\rank(A)=m$.
\end{proposition}

\begin{proof}
Writing $A=[c_1\ \cdots\ c_n]$ and $x=(x_1,\ldots,x_n)^\top$ gives
\[
Ax=x_1c_1+\cdots+x_nc_n.
\]
Hence every vector in $\im(A)$ is in the span of the columns, and every linear
combination of columns is equal to $Ax$ for the vector of coefficients $x$.
Surjectivity means this column span is all of $\R^m$, which is equivalent to
rank $m$.
\end{proof}

\begin{definition}[Kernel and injectivity]
The kernel, or null space, of $A$ is
\[
\ker(A)=\{x\in\R^n:Ax=0\}.
\]
It is the set of directions collapsed by $A$.

The map $T_A:\R^n\to\R^m$ is injective if different inputs give different
outputs.
\end{definition}

\begin{proposition}[Kernel criterion for injectivity]
For a linear map $T_A:\R^n\to\R^m$, injectivity is equivalent to
\[
\ker(A)=\{0\}.
\]
Equivalently, $A$ has rank $n$.
\end{proposition}

\begin{proof}
If $A$ is injective and $Ax=0=A0$, then $x=0$, so $\ker(A)=\{0\}$. Conversely,
if $\ker(A)=\{0\}$ and $Ax=Ay$, then $A(x-y)=0$, hence $x-y=0$ and $x=y$.
By rank-nullity, $\ker(A)=\{0\}$ is equivalent to $\rank(A)=n$.
\end{proof}

\begin{definition}[Rank]
The rank of $A$ is
\[
\rank(A)=\dim \im(A).
\]
\end{definition}

\begin{proposition}[Rank-nullity theorem]
The rank-nullity theorem says that for $A\in\M_{m,n}(\R)$,
\[
\rank(A)+\dim\ker(A)=n,
\]
where $n$ is the number of columns of $A$.
\end{proposition}

\begin{proof}
Row-reduce $A$ to echelon form. Row operations amount to multiplying on the
left by an invertible matrix, so they preserve the kernel and the rank. If the
echelon form has $r$ pivot columns, then $r=\rank(A)$ and the remaining $n-r$
variables are free. Therefore $\dim\ker(A)=n-r$, giving
$\rank(A)+\dim\ker(A)=r+(n-r)=n$.
\end{proof}

\begin{proposition}[Equivalent conditions for square matrices]
For $A\in\M_n(\R)$, injectivity, surjectivity, and invertibility are equivalent:
\[
A\text{ injective}
\quad\Longleftrightarrow\quad
A\text{ surjective}
\quad\Longleftrightarrow\quad
A\text{ invertible}.
\]
Equivalently:
\begin{itemize}
\item $\ker(A)=\{0\}$;
\item $\im(A)=\R^n$;
\item $\rank(A)=n$;
\item $\det(A)\neq0$;
\item for every $b\in\R^n$, $Ax=b$ has a unique solution.
\end{itemize}
\end{proposition}

\begin{proof}
For a square $n\times n$ matrix, rank-nullity shows that
$\ker(A)=\{0\}$ is equivalent to $\rank(A)=n$. Having rank $n$ means the columns
span $\R^n$, so $\im(A)=\R^n$. Thus injectivity, surjectivity, and full rank are
equivalent. Full rank is equivalent to $\det(A)\neq0$, hence to invertibility.
If $A$ is invertible, $Ax=b$ has the unique solution $A^{-1}b$; conversely, if
every $b$ has a unique solution, the columns form a basis and $A$ is invertible.
\end{proof}

\begin{proposition}[Solution set of a linear system]
For a general matrix $A\in\M_{m,n}(\R)$:
\begin{itemize}
\item $Ax=b$ has at least one solution if and only if $b\in\im(A)$.
\item If $x_0$ is one solution, then all solutions are
\[
x=x_0+z,
\qquad z\in\ker(A).
\]
\item If $\ker(A)=\{0\}$, then a solution, if it exists, is unique.
\end{itemize}
\end{proposition}

\begin{proof}
The first statement is exactly the definition of $\im(A)$. If $x_0$ solves
$Ax=b$, then $A(x_0+z)=b$ for every $z\in\ker(A)$. Conversely, if $x$ is any
solution, then $A(x-x_0)=0$, so $x-x_0\in\ker(A)$ and $x=x_0+z$. If the kernel
is $\{0\}$, this representation has no nonzero direction to add, so the
solution is unique.
\end{proof}

\begin{definition}[Orthogonal vectors and orthogonal complement]
Two vectors $x,y\in\R^m$ are orthogonal, written $x\perp y$, if
\[
x^\top y=0.
\]
If $V\subset\R^m$ is a subspace, its orthogonal complement is
\[
V^\perp
=\{y\in\R^m:y^\top v=0\text{ for every }v\in V\}.
\]
Thus $V^\perp$ contains exactly the directions perpendicular to the whole
subspace $V$.
\end{definition}

\begin{proposition}[Fundamental orthogonality relation]
For $A\in\M_{m,n}(\R)$,
\[
\im(A)^\perp=\ker(A^\top),
\qquad
\im(A)=(\ker A^\top)^\perp.
\]
\end{proposition}

\begin{proof}
For $y\in\R^m$, the condition $y\in\im(A)^\perp$ means
$y^\top Ax=0$ for every $x\in\R^n$. This is equivalent to
$(A^\top y)^\top x=0$ for every $x$, hence to $A^\top y=0$. This proves
$\im(A)^\perp=\ker(A^\top)$. The second identity follows by taking orthogonal
complements in finite dimension.
\end{proof}

\subsection*{Worked example: kernel and image}

Let
\[
A=
\begin{pmatrix}
1&-1&0\\
-1&2&-1\\
0&-1&1
\end{pmatrix}.
\]
To find the kernel, solve $Ax=0$:
\[
\begin{cases}
x_1-x_2=0,\\
-x_1+2x_2-x_3=0,\\
-x_2+x_3=0.
\end{cases}
\]
The first and third equations give $x_1=x_2$ and $x_3=x_2$. Hence
\[
x=t\begin{pmatrix}1\\1\\1\end{pmatrix},
\qquad
\ker(A)=\Span\left\{\begin{pmatrix}1\\1\\1\end{pmatrix}\right\}.
\]
Since the kernel is one-dimensional and $A$ has three columns,
\[
\rank(A)=3-1=2.
\]
The image is a two-dimensional plane in $\R^3$.

Because $A$ is symmetric, $A^\top=A$. Using the fundamental relation
\[
\im(A)=(\ker A^\top)^\perp,
\]
we get
\[
\im(A)=(\ker A)^\perp
=\left\{y\in\R^3:y_1+y_2+y_3=0\right\}.
\]

\subsection*{Numerical example: a rectangular matrix}

Let
\[
B=
\begin{pmatrix}
1&2&3\\
2&4&6
\end{pmatrix}.
\]
The second row is twice the first, so the two rows are dependent. The columns are
\[
c_1=\begin{pmatrix}1\\2\end{pmatrix},
\quad
c_2=\begin{pmatrix}2\\4\end{pmatrix}=2c_1,
\quad
c_3=\begin{pmatrix}3\\6\end{pmatrix}=3c_1.
\]
Thus
\[
\im(B)=\Span\left\{\begin{pmatrix}1\\2\end{pmatrix}\right\},
\qquad
\rank(B)=1.
\]
By rank-nullity,
\[
\dim\ker(B)=3-\rank(B)=2.
\]
Indeed, $Bx=0$ means
\[
x_1+2x_2+3x_3=0.
\]
Taking $x_2=s$ and $x_3=t$ gives $x_1=-2s-3t$, so
\[
\ker(B)=
\Span\left\{
\begin{pmatrix}-2\\1\\0\end{pmatrix},
\begin{pmatrix}-3\\0\\1\end{pmatrix}
\right\}.
\]
Here $B:\R^3\to\R^2$ is neither injective nor surjective: it has a nontrivial
kernel, and its image is only a line in $\R^2$.

\section{Classical Decompositions}

\subsection*{Why decompose matrices?}

Matrix decompositions rewrite a matrix in a form that makes its action easier to
understand or compute. A decomposition often answers one of the following
questions:
\begin{itemize}
\item Can we solve a system faster?
\item Can we understand the main directions of a transformation?
\item Can we separate scaling, rotation, and correlation effects?
\item Can we simulate or transform correlated random variables?
\end{itemize}

\subsection{Eigenvalue Diagonalization}

\subsubsection{Definition and intuition}

The simplest possible matrix is a diagonal matrix:
\[
D=\diag(\lambda_1,\ldots,\lambda_n).
\]
Its action is component-by-component:
\[
Dx=
\begin{pmatrix}
\lambda_1x_1\\
\vdots\\
\lambda_nx_n
\end{pmatrix}.
\]
Diagonalization asks whether a general matrix can become diagonal after changing
coordinates.

\begin{definition}[Diagonalizable matrix]
A matrix $A\in\M_n(\R)$ is diagonalizable over $\R$ if there
exist an invertible matrix $P$ and a diagonal matrix $D$ such that
\[
A=PDP^{-1}.
\]
Equivalently, in the coordinate system defined by $P$, the matrix $A$ acts like
the diagonal matrix $D$.
\end{definition}

\begin{definition}[Eigenvalue, eigenvector, eigenspace]
An eigenvector of $A\in\M_n(\R)$ is a nonzero vector $v$ such that
\[
Av=\lambda v.
\]
The scalar $\lambda$ is the associated eigenvalue. The direction of $v$ is
preserved by $A$; only its magnitude and possibly its orientation change.
The eigenspace associated with $\lambda$ is
\[
E_\lambda=\ker(A-\lambda I).
\]
\end{definition}

\begin{proposition}[Eigenbasis criterion]
If there is a basis of eigenvectors $(v_1,\ldots,v_n)$, then $A$ is
diagonalizable. Put these eigenvectors in the columns of $P$:
\[
P=
\begin{pmatrix}
|&&|\\
v_1&\cdots&v_n\\
|&&|
\end{pmatrix},
\qquad
D=\diag(\lambda_1,\ldots,\lambda_n).
\]
Then
\[
AP=PD,
\qquad\text{hence}\qquad
A=PDP^{-1}.
\]
Conversely, if $A=PDP^{-1}$ with $D$ diagonal, the columns of $P$ form a basis
of eigenvectors of $A$.
\end{proposition}

\begin{proof}
If $Av_i=\lambda_i v_i$ for each basis vector, then the $i$th column of $AP$ is
$Av_i$ and the $i$th column of $PD$ is $\lambda_i v_i$. Hence $AP=PD$, and
since $P$ is invertible, $A=PDP^{-1}$. Conversely, if $A=PDP^{-1}$, then
$AP=PD$. The $i$th column gives $Ap_i=\lambda_i p_i$, where $p_i$ is the $i$th
column of $P$. Since $P$ is invertible, these columns form a basis.
\end{proof}

\begin{definition}[Characteristic polynomial]
Eigenvalues are found by solving
\[
\det(A-\lambda I)=0.
\]
The polynomial
\[
p_A(\lambda)=\det(A-\lambda I)
\]
is the characteristic polynomial.
\end{definition}

\begin{proposition}[Eigenvalues are roots of the characteristic polynomial]
A real number $\lambda$ is an eigenvalue of $A$ if and only if
$p_A(\lambda)=0$.
\end{proposition}

\begin{proof}
The equation $Av=\lambda v$ with $v\neq0$ is equivalent to
$(A-\lambda I)v=0$ having a nonzero solution. This happens if and only if
$A-\lambda I$ is not injective, equivalently not invertible. For a square
matrix, non-invertibility is equivalent to determinant zero.
\end{proof}

\subsubsection{Computing eigenvalues and eigenvectors}

Let
\[
A=
\begin{pmatrix}
4&1\\
2&3
\end{pmatrix}.
\]
Then
\[
A-\lambda I=
\begin{pmatrix}
4-\lambda&1\\
2&3-\lambda
\end{pmatrix}
\]
and
\[
\det(A-\lambda I)
=(4-\lambda)(3-\lambda)-2
=\lambda^2-7\lambda+10
=(\lambda-5)(\lambda-2).
\]
Thus $\lambda_1=5$ and $\lambda_2=2$.

For $\lambda_1=5$,
\[
A-5I=
\begin{pmatrix}
-1&1\\
2&-2
\end{pmatrix},
\]
so $-v_1+v_2=0$. We may choose
\[
u_1=\begin{pmatrix}1\\1\end{pmatrix}.
\]
For $\lambda_2=2$,
\[
A-2I=
\begin{pmatrix}
2&1\\
2&1
\end{pmatrix},
\]
so $2v_1+v_2=0$. We may choose
\[
u_2=\begin{pmatrix}1\\-2\end{pmatrix}.
\]

\subsubsection{Computational purpose}

Diagonalization turns matrix operations into scalar operations. If $A=PDP^{-1}$,
then
\[
A^k=PD^kP^{-1},
\qquad
D^k=\diag(\lambda_1^k,\ldots,\lambda_n^k).
\]
This is useful for powers of matrices, linear recurrences, Markov chains,
discrete-time dynamical systems, covariance analysis, and PCA.

\subsubsection{Algorithm}

\begin{enumerate}
\item Compute $p_A(\lambda)=\det(A-\lambda I)$.
\item Find the eigenvalues.
\item For each eigenvalue $\lambda$, solve $(A-\lambda I)v=0$.
\item Check whether the eigenvectors found form a basis of $\R^n$.
\item If yes, build $P$ from eigenvectors and $D$ from eigenvalues.
\end{enumerate}

\begin{proposition}[Distinct real eigenvalues imply diagonalizability]
If $A\in\M_n(\R)$ has $n$ distinct real eigenvalues, then $A$ is diagonalizable
over $\R$.
\end{proposition}

\begin{proof}
Choose one eigenvector $v_i$ for each distinct eigenvalue $\lambda_i$. Suppose
$\sum_i \alpha_i v_i=0$. For a fixed $j$, apply the polynomial
$q_j(A)=\prod_{i\neq j}(A-\lambda_i I)$ to this relation. Every term vanishes
except the $j$th one, giving
\[
\alpha_j\prod_{i\neq j}(\lambda_j-\lambda_i)v_j=0.
\]
The product is nonzero and $v_j\neq0$, so $\alpha_j=0$. Thus the $v_i$ are
linearly independent. There are $n$ of them, so they form an eigenbasis.
\end{proof}

\begin{definition}[Algebraic and geometric multiplicities]
The algebraic multiplicity of an eigenvalue is its multiplicity as a root of
the characteristic polynomial. The geometric multiplicity is the dimension of
the eigenspace
\[
E_\lambda=\ker(A-\lambda I).
\]
\end{definition}

\begin{proposition}[Multiplicity criterion]
Always,
\[
1\le \dim E_\lambda\le \text{algebraic multiplicity of }\lambda.
\]
The matrix is diagonalizable if and only if the sum of the dimensions of all
eigenspaces is $n$. Equivalently, when the characteristic polynomial splits
over $\R$, $A$ is diagonalizable if and only if for every eigenvalue the
geometric multiplicity equals the algebraic multiplicity.
\end{proposition}

\begin{proof}
The eigenspace is nonzero by definition of eigenvalue, so its dimension is at
least one. If $g=\dim E_\lambda$, choose a basis whose first $g$ vectors span
$E_\lambda$. In this basis, $A$ has a block form with $\lambda I_g$ in the
upper-left block, so the characteristic polynomial contains the factor
$(\lambda-t)^g$ up to sign. Hence $g$ is at most the algebraic multiplicity.
Eigenspaces associated with distinct eigenvalues are linearly independent, as
proved above. Therefore $A$ is diagonalizable exactly when the eigenspaces
together provide $n$ independent eigenvectors, i.e. when their dimensions sum
to $n$.
\end{proof}

\begin{proposition}[Immediate sufficient cases]
Every diagonal matrix is diagonalizable. Every triangular matrix with distinct
diagonal entries is diagonalizable.
\end{proposition}

\begin{proof}
For a diagonal matrix, the canonical basis vectors are eigenvectors. For a
triangular matrix, the eigenvalues are the diagonal entries because
$A-\lambda I$ is triangular and its determinant is the product of diagonal
entries. If these diagonal entries are distinct, the previous proposition
applies.
\end{proof}

\subsubsection{Worked example: a general two-by-two matrix}

Use the previous matrix:
\[
A=
\begin{pmatrix}
4&1\\
2&3
\end{pmatrix}.
\]
We found eigenvalues $\lambda_1=5$, $\lambda_2=2$, with eigenvectors
\[
u_1=\begin{pmatrix}1\\1\end{pmatrix},
\qquad
u_2=\begin{pmatrix}1\\-2\end{pmatrix}.
\]
Set
\[
P=
\begin{pmatrix}
1&1\\
1&-2
\end{pmatrix},
\qquad
D=
\begin{pmatrix}
5&0\\
0&2
\end{pmatrix}.
\]
Since
\[
\det(P)=1(-2)-1(1)=-3\neq0,
\]
$P$ is invertible. Hence
\[
A=PDP^{-1}.
\]
Here
\[
P^{-1}
=
\frac{1}{-3}
\begin{pmatrix}
-2&-1\\
-1&1
\end{pmatrix}
=
\begin{pmatrix}
2/3&1/3\\
1/3&-1/3
\end{pmatrix}.
\]
This gives an efficient formula for powers:
\[
A^k
=
P
\begin{pmatrix}
5^k&0\\
0&2^k
\end{pmatrix}
P^{-1}.
\]

\subsubsection{Worked example: a triangular matrix}

Let
\[
A=
\begin{pmatrix}
3&1\\
0&2
\end{pmatrix}.
\]
Since $A$ is triangular, its eigenvalues are the diagonal entries:
\[
\lambda_1=3,
\qquad
\lambda_2=2.
\]
For $\lambda=3$,
\[
A-3I=
\begin{pmatrix}
0&1\\
0&-1
\end{pmatrix},
\]
so $v_2=0$ and we may choose
\[
u_1=\begin{pmatrix}1\\0\end{pmatrix}.
\]
For $\lambda=2$,
\[
A-2I=
\begin{pmatrix}
1&1\\
0&0
\end{pmatrix},
\]
so $v_1+v_2=0$ and we may choose
\[
u_2=\begin{pmatrix}-1\\1\end{pmatrix}.
\]
Thus
\[
P=\begin{pmatrix}1&-1\\0&1\end{pmatrix},
\qquad
D=\begin{pmatrix}3&0\\0&2\end{pmatrix},
\qquad
A=PDP^{-1}.
\]

\subsubsection{Worked example: a covariance-style matrix}

Let
\[
C=
\begin{pmatrix}
1&0.8\\
0.8&1
\end{pmatrix}.
\]
The characteristic polynomial is
\[
\det(C-\lambda I)
=
(1-\lambda)^2-0.8^2.
\]
Thus
\[
1-\lambda=\pm0.8,
\]
so
\[
\lambda_1=1.8,
\qquad
\lambda_2=0.2.
\]
For $\lambda_1=1.8$, an eigenvector is $(1,1)^\top$. For $\lambda_2=0.2$, an
eigenvector is $(1,-1)^\top$. Normalize them:
\[
q_1=\frac{1}{\sqrt2}\begin{pmatrix}1\\1\end{pmatrix},
\qquad
q_2=\frac{1}{\sqrt2}\begin{pmatrix}1\\-1\end{pmatrix}.
\]
Set
\[
Q=
\frac{1}{\sqrt2}
\begin{pmatrix}
1&1\\
1&-1
\end{pmatrix},
\qquad
D=
\begin{pmatrix}
1.8&0\\
0&0.2
\end{pmatrix}.
\]
Since $Q$ is orthogonal, $Q^{-1}=Q^\top$. Therefore
\[
C=QDQ^\top.
\]
Interpretation: most variance lies in the common direction $(1,1)$, while much
less variance lies in the spread direction $(1,-1)$.

\subsubsection{Counterexample: a non-diagonalizable matrix}

Let
\[
J=
\begin{pmatrix}
1&1\\
0&1
\end{pmatrix}.
\]
The only eigenvalue is $\lambda=1$, with algebraic multiplicity $2$. But
\[
J-I=
\begin{pmatrix}
0&1\\
0&0
\end{pmatrix},
\]
so
\[
(J-I)v=0
\quad\Longrightarrow\quad
v_2=0.
\]
The eigenspace is
\[
E_1=\Span\left\{\begin{pmatrix}1\\0\end{pmatrix}\right\},
\]
which has dimension $1$, not $2$. Therefore $J$ is not diagonalizable.

\subsubsection{Diagonalization over the reals versus the complex numbers}

The rotation matrix
\[
R=
\begin{pmatrix}
0&-1\\
1&0
\end{pmatrix}
\]
has characteristic polynomial
\[
\det(R-\lambda I)=\lambda^2+1.
\]
It has no real eigenvalues, so it is not diagonalizable over $\R$. It is
diagonalizable over $\mathbb{C}$, with eigenvalues $i$ and $-i$. This distinction
matters whenever we specify the field over which diagonalization is being
performed.

\subsubsection{Application: a discrete-time dynamical system}

Consider the recursion
\[
x_{k+1}=Ax_k,
\qquad
A=\begin{pmatrix}3&1\\0&2\end{pmatrix},
\qquad
x_0=\begin{pmatrix}0\\1\end{pmatrix}.
\]
The triangular matrix used in the slides was diagonalized as
\[
A=PDP^{-1},\qquad
P=\begin{pmatrix}1&-1\\0&1\end{pmatrix},
\qquad
D=\begin{pmatrix}3&0\\0&2\end{pmatrix}.
\]
Since
\[
P^{-1}=\begin{pmatrix}1&1\\0&1\end{pmatrix},
\]
we obtain, for every integer $k\geq0$,
\[
A^k=PD^kP^{-1}
=\begin{pmatrix}
3^k&3^k-2^k\\
0&2^k
\end{pmatrix}.
\]
Therefore the complete trajectory is
\[
x_k=A^kx_0
=\begin{pmatrix}3^k-2^k\\2^k\end{pmatrix}.
\]
Instead of iterating a coupled two-dimensional recursion, diagonalization
separates the system into the two scalar growth modes $3^k$ and $2^k$. The
mode associated with the largest eigenvalue eventually dominates.

\subsection{Spectral Decomposition}

\begin{proposition}[Spectral theorem]
If $A$ is real and symmetric, then it can be diagonalized using an orthogonal
matrix:
\[
A=QDQ^\top.
\]
The eigenvalues are real, and eigenvectors associated with distinct eigenvalues
are orthogonal.
\end{proposition}

\begin{proof}
The Rayleigh quotient $x^\top Ax$ on the unit sphere attains a maximum because
the unit sphere is compact. At a maximizer $q_1$, the Lagrange multiplier
condition gives $Aq_1=\lambda_1 q_1$, so $q_1$ is an eigenvector. If
$y\perp q_1$, then $q_1^\top Ay=(Aq_1)^\top y=\lambda_1 q_1^\top y=0$, so the
orthogonal complement of $q_1$ is invariant under $A$. Applying the same
argument by induction gives an orthonormal eigenbasis. Placing these vectors in
the columns of $Q$ gives $Q^\top Q=I$ and $A=QDQ^\top$.
\end{proof}

This theorem is central for covariance matrices, PCA, and risk-factor
decompositions.

\subsubsection{Application: common and spread risk factors}

Return to the covariance matrix
\[
C=\begin{pmatrix}1&0.8\\0.8&1\end{pmatrix}.
\]
Its orthonormal eigenvectors and eigenvalues are
\[
q_1=\frac{1}{\sqrt2}\begin{pmatrix}1\\1\end{pmatrix},
\quad \lambda_1=1.8,
\qquad
q_2=\frac{1}{\sqrt2}\begin{pmatrix}1\\-1\end{pmatrix},
\quad \lambda_2=0.2.
\]
For any unit portfolio direction $q$, its return variance is $q^\top Cq$.
Consequently,
\[
q_1^\top Cq_1=1.8,
\qquad
q_2^\top Cq_2=0.2.
\]
The first portfolio moves both assets in the same direction and represents a
common factor. The second is a long--short spread. The spectral decomposition
therefore produces orthogonal portfolios and measures exactly how much risk is
carried by each one.

\subsection{Cholesky Decomposition}

\begin{proposition}[Cholesky decomposition]
If $A$ is symmetric positive definite, then there exists a unique lower
triangular matrix $L$ with positive diagonal entries such that
\[
A=LL^\top.
\]
This is the Cholesky decomposition.
\end{proposition}

\begin{proof}
Proceed by induction on the dimension. Write
\[
A=\begin{pmatrix}a&r^\top\\ r&B\end{pmatrix},
\qquad a>0.
\]
Set $\ell_{11}=\sqrt a$ and $\ell=r/\sqrt a$. The Schur complement
$S=B-\ell\ell^\top=B-rr^\top/a$ is positive definite: for $y\neq0$, choose
$\alpha=-r^\top y/a$ and use positive definiteness of
$(\alpha,y)^\top A(\alpha,y)$. By induction, $S=MM^\top$ with $M$ lower
triangular and positive diagonal. Then
\[
L=\begin{pmatrix}\sqrt a&0\\ \ell&M\end{pmatrix}
\]
satisfies $A=LL^\top$. Uniqueness follows by matching entries from the first
column onward; the positive diagonal fixes the sign at each step.
\end{proof}

\subsubsection{Worked Cholesky factorization}

Let
\[
A=
\begin{pmatrix}
4&2\\
2&3
\end{pmatrix}.
\]
Seek
\[
L=
\begin{pmatrix}
\ell_{11}&0\\
\ell_{21}&\ell_{22}
\end{pmatrix}
\]
such that $A=LL^\top$. Then
\[
LL^\top=
\begin{pmatrix}
\ell_{11}^2&\ell_{11}\ell_{21}\\
\ell_{11}\ell_{21}&\ell_{21}^2+\ell_{22}^2
\end{pmatrix}.
\]
Matching entries gives
\[
\ell_{11}=2,
\qquad
2\ell_{21}=2\Rightarrow \ell_{21}=1,
\]
and
\[
\ell_{21}^2+\ell_{22}^2=3
\quad\Rightarrow\quad
1+\ell_{22}^2=3
\quad\Rightarrow\quad
\ell_{22}=\sqrt2.
\]
Therefore
\[
A=
\begin{pmatrix}
2&0\\
1&\sqrt2
\end{pmatrix}
\begin{pmatrix}
2&1\\
0&\sqrt2
\end{pmatrix}.
\]

\subsubsection{Why Cholesky matters}

Cholesky is widely used to:
\begin{itemize}
\item solve positive definite linear systems efficiently;
\item simulate correlated Gaussian vectors;
\item work with covariance matrices;
\item implement numerical optimization algorithms.
\end{itemize}
If $Z\sim N(0,I)$ and $\Sigma=LL^\top$, then
\[
X=LZ
\]
has covariance matrix $\Sigma$.

\subsubsection{Application: simulation of correlated returns}

Let $Z_1,Z_2$ be independent standard Gaussian variables and take the factor
computed above,
\[
L=\begin{pmatrix}2&0\\1&\sqrt2\end{pmatrix}.
\]
Define
\[
X=L\begin{pmatrix}Z_1\\Z_2\end{pmatrix}
=\begin{pmatrix}2Z_1\\Z_1+\sqrt2 Z_2\end{pmatrix}.
\]
Since $\operatorname{Cov}(Z)=I_2$,
\[
\operatorname{Cov}(X)=L\operatorname{Cov}(Z)L^\top
=LL^\top
=\begin{pmatrix}4&2\\2&3\end{pmatrix}.
\]
Thus $X_1$ and $X_2$ have variances $4$ and $3$, covariance $2$, and
correlation
\[
\rho_{12}=\frac{2}{\sqrt{4}\sqrt{3}}=\frac{1}{\sqrt3}.
\]
This gives a practical simulation recipe: generate independent standard
Gaussian shocks and multiply them by $L$ to impose the desired covariance
structure.

\section{Statistics and Finance Applications}

\begin{definition}[Empirical covariance matrix]
If $X\in\R^{n\times d}$ is a centered data matrix, the empirical covariance
matrix is
\[
\widehat{\Sigma}=\frac{1}{n-1}X^\top X.
\]
\end{definition}

\begin{proposition}[Covariance matrices are positive semidefinite]
The matrix $\widehat{\Sigma}$ is symmetric and positive semidefinite.
\end{proposition}

\begin{proof}
Symmetry follows from $(X^\top X)^\top=X^\top X$. For every $z\in\R^d$,
\[
z^\top\widehat{\Sigma}z
=\frac{1}{n-1}z^\top X^\top Xz
=\frac{1}{n-1}\|Xz\|^2\ge0.
\]
\end{proof}

\subsection*{Numerical covariance example}

Let
\[
X=
\begin{pmatrix}
0.01&-0.03\\
0.02&0.01\\
-0.01&0
\end{pmatrix}.
\]
Then
\[
X^\top X
=
\begin{pmatrix}
0.0006&-0.0001\\
-0.0001&0.0010
\end{pmatrix}.
\]

\begin{definition}[Portfolio return]
For asset returns $r$ with covariance matrix $\Sigma$ and portfolio weights $w$,
the portfolio return is
\[
R_p=w^\top r,
\]
where the weights $w_i$ represent the portfolio allocation.
\end{definition}

\begin{proposition}[Portfolio variance formula]
The variance of the portfolio return is
\[
\Var(R_p)=w^\top\Sigma w.
\]
\end{proposition}

\begin{proof}
Let $\mu=\mathbb{E}[r]$ and $\Sigma=\mathbb{E}[(r-\mu)(r-\mu)^\top]$. Then
\[
\Var(w^\top r)
=\mathbb{E}\left[(w^\top(r-\mu))^2\right]
=w^\top\mathbb{E}\left[(r-\mu)(r-\mu)^\top\right]w
=w^\top\Sigma w.
\]
\end{proof}

For
\[
\Sigma=
\begin{pmatrix}
0.04&0.018\\
0.018&0.09
\end{pmatrix},
\qquad
w=\begin{pmatrix}0.6\\0.4\end{pmatrix},
\]
we compute
\[
w^\top\Sigma w
=0.6^2(0.04)+2(0.6)(0.4)(0.018)+0.4^2(0.09)
=0.03744.
\]
This quadratic form decomposes risk into individual variance terms and a
covariance term.

\begin{proposition}[PCA as covariance diagonalization]
For centered data, PCA diagonalizes the covariance matrix
\[
\widehat{\Sigma}=\frac{1}{n-1}X^\top X.
\]
The principal directions are eigenvectors of $\widehat{\Sigma}$, and the
corresponding eigenvalues are variances along those directions.
\end{proposition}

\begin{proof}
By the spectral theorem, $\widehat{\Sigma}=QDQ^\top$ with $Q$ orthogonal and
$D$ diagonal. If $q_i$ is a column of $Q$, then the variance of the projected
data $Xq_i$ is
\[
\frac{1}{n-1}\|Xq_i\|^2
=q_i^\top\widehat{\Sigma}q_i
=q_i^\top QDQ^\top q_i
=D_{ii}.
\]
Thus eigenvectors give the orthogonal directions of variation and eigenvalues
give the amount of variance in those directions.
\end{proof}

Let
\[
X=
\begin{pmatrix}
-1&-1\\
0&0\\
1&1
\end{pmatrix}.
\]
Then
\[
\widehat{\Sigma}=\frac{1}{2}X^\top X
=
\begin{pmatrix}1&1\\1&1\end{pmatrix}.
\]
The eigenvalues are $2$ and $0$, with eigenvectors
\[
q_1=\frac{1}{\sqrt 2}\begin{pmatrix}1\\1\end{pmatrix},
\qquad
q_2=\frac{1}{\sqrt 2}\begin{pmatrix}1\\-1\end{pmatrix}.
\]
All the variance lies along the line $x_1=x_2$.

\subsection*{Projection onto the first principal component}

The first principal direction is $q_1=(1,1)^\top/\sqrt2$. The projected
coordinates are
\[
Xq_1=
\begin{pmatrix}
-1&-1\\
0&0\\
1&1
\end{pmatrix}
\frac{1}{\sqrt2}
\begin{pmatrix}
1\\1
\end{pmatrix}
=
\begin{pmatrix}
-\sqrt2\\
0\\
\sqrt2
\end{pmatrix}.
\]
This keeps the direction where the data vary. Projection onto $q_2$ would give
zero for every observation, because there is no variation in the orthogonal
direction.

\section{Markov Chains and Transition Matrices}

\subsection{From conditional probabilities to a matrix}

Let $(X_n)_{n\geq0}$ take values in a finite state space
$\{1,\ldots,m\}$. It is a Markov chain if the conditional distribution of the
next state depends on the past only through the current state:
\[
\mathbb{P}(X_{n+1}=i\mid X_n=j,X_{n-1},\ldots,X_0)
=\mathbb{P}(X_{n+1}=i\mid X_n=j).
\]

\begin{definition}[Transition matrix]
Using the column-vector convention of these notes, define
\[
P_{ij}=\mathbb{P}(X_{n+1}=i\mid X_n=j).
\]
Thus column $j$ contains the probabilities of the next state given that the
current state is $j$. Every entry is nonnegative and every column sums to one:
\[
P_{ij}\geq0,
\qquad
\sum_{i=1}^mP_{ij}=1.
\]
Such a matrix is called column-stochastic.
\end{definition}

Some books instead use row probability vectors and row-stochastic matrices.
Both conventions are equivalent after transposition; the important point is to
state the convention and use it consistently.

\subsection{Distribution dynamics and powers of the matrix}

Let
\[
\mu_n=
\begin{pmatrix}
\mathbb{P}(X_n=1)\\
\vdots\\
\mathbb{P}(X_n=m)
\end{pmatrix}.
\]
By the law of total probability,
\[
(\mu_{n+1})_i
=\sum_{j=1}^m
\mathbb{P}(X_{n+1}=i\mid X_n=j)\mathbb{P}(X_n=j)
=\sum_{j=1}^mP_{ij}(\mu_n)_j.
\]
Therefore
\[
\boxed{\mu_{n+1}=P\mu_n},
\qquad
\boxed{\mu_n=P^n\mu_0}.
\]
More generally, the entry $(P^n)_{ij}$ is the probability of being in state
$i$ after $n$ steps when the chain starts from state $j$.

This is precisely why matrix powers matter in probability: multiplying
transition matrices adds time steps.

\subsection{A two-state credit example}

Consider two credit states, Good (G) and Bad (B). From G, the rating stays Good
with probability $0.9$ and deteriorates with probability $0.1$. From B, it
recovers with probability $0.4$ and stays Bad with probability $0.6$. With
columns indexed by the current state,
\[
P=
\begin{pmatrix}
0.9&0.4\\
0.1&0.6
\end{pmatrix}.
\]
If the obligor starts in state G, then $\mu_0=(1,0)^\top$ and
\[
\mu_1=P\mu_0=
\begin{pmatrix}0.9\\0.1\end{pmatrix},
\qquad
\mu_2=P^2\mu_0=
\begin{pmatrix}0.85\\0.15\end{pmatrix}.
\]
Thus the two-year probability of being in the Bad state is $0.15$.

\subsection{Stationary distributions and diagonalization}

\begin{definition}[Stationary distribution]
A probability vector $\pi$ is stationary if
\[
P\pi=\pi,
\qquad
\mathbf{1}^\top\pi=1,
\qquad
\pi_i\geq0.
\]
It is therefore an eigenvector of $P$ associated with eigenvalue $1$, normalized
so that its entries sum to one.
\end{definition}

For the credit matrix,
\[
P\pi=\pi
\quad\Longleftrightarrow\quad
0.1\pi_G=0.4\pi_B,
\qquad
\pi_G+\pi_B=1,
\]
so
\[
\pi=\begin{pmatrix}0.8\\0.2\end{pmatrix}.
\]
The eigenvalues of $P$ are $1$ and $0.5$. Since the matrix is diagonalizable,
its powers separate the persistent stationary mode from the transient mode. In
fact, when $\mu_0=(1,0)^\top$,
\[
\mu_n
=\begin{pmatrix}0.8\\0.2\end{pmatrix}
+0.2(0.5)^n\begin{pmatrix}1\\-1\end{pmatrix}.
\]
Because $(0.5)^n\to0$, the distribution converges to $\pi$. This illustrates a
central link between linear algebra and probability:
\begin{quote}
eigenvalue $1$ determines the long-run distribution, while eigenvalues of
modulus less than one determine the speed at which transient effects disappear.
\end{quote}

\end{document}
